% general-case-algorithm.tex — \input'd from algorithms.tex
% Conventions: thesis/CLAUDE.md
%
% Sources: Jörn's dictation + HK2017 paper. Notation per correspondence.tex.
%
% Structure: Main Result theorem and algorithm, Well-definedness lemma

% QC: polytope per Definition~\ref{def:polytope}, facet data per Definition~\ref{def:facets},
% capacity per Definition~\ref{def:ehz-capacity}.

% Jörn: text approved (06a976c) — from \begin{theorem} to \end{theorem}
\begin{theorem}[EHZ capacity formula {\cite{HK2017}}]\label{thm:main}
Let $K \subset \mathbb{R}^4$ be a polytope
with $F$ facets, outward unit normals $n_1, \ldots, n_F \in S^3$,
and heights $h_1, \ldots, h_F > 0$.
Then
\begin{equation}\label{eq:ehz-main}
  c_{\mathrm{EHZ}}(K) = \frac{1}{2}
  \left[
    \max_{\substack{
      \sigma \in S_F \\
      \beta \in M(K)
    }}
    \sum_{1 \leq j < i \leq F}
      \beta_{\sigma(i)} \beta_{\sigma(j)} \,
      \omega_0(n_{\sigma(j)}, n_{\sigma(i)})
  \right]^{-1},
\end{equation}
where
\begin{equation}\label{eq:constraint-set}
  M(K) = \left\{
    (\beta_i)_{i=1}^{F} :
    \beta_i \geq 0, \quad  % not >: M(K) includes boundary where some \beta_i = 0
    \sum_{i=1}^{F} \beta_i h_i = 1, \quad
    \sum_{i=1}^{F} \beta_i n_i = 0
  \right\}.
\end{equation}
\end{theorem}

% Proof: Section~\ref{sec:simple-minimizer} (search space)
% and Section~\ref{sec:algorithm-correctness} (algorithm correctness).

% Jörn: text approved (06a976c) — from \begin{algorithm} to \end{algorithm}
\begin{algorithm}[EHZ capacity]\label{alg:ehz}
Computes~\eqref{eq:ehz-main}.

\smallskip\noindent
\emph{Input:} $(n_1, h_1), \ldots, (n_F, h_F)$.
\quad
\emph{Output:} $c_{\mathrm{EHZ}}(K)$.

\smallskip\noindent
For each nonempty subset $S \subseteq \{1, \ldots, F\}$
and each permutation $\sigma$ of $S$
(identified up to cyclic shifts):

\begin{enumerate}
\item \textbf{Define}
  the normal matrix $N \in \mathbb{R}^{|S| \times 4}$,
  action matrix $H \in \mathbb{R}^{|S| \times |S|}$,
  and height vector $\eta \in \mathbb{R}^{|S|}$:
  \begin{align*}
    N_{id} &= n_{\sigma(i),d},
      & i &= 1, \ldots, |S|, \quad d = 1, \ldots, 4, \\
    H_{ij} &= \begin{cases}
      \omega_0(n_{\sigma(i)}, n_{\sigma(j)}) & \text{if } i < j, \\
      0 & \text{if } i = j, \\
      \omega_0(n_{\sigma(j)}, n_{\sigma(i)}) & \text{if } i > j,
    \end{cases} \\
    \eta_i &= h_{\sigma(i)},
      & i &= 1, \ldots, |S|.
  \end{align*}

\item \textbf{Solve.}
  \begin{equation}\label{eq:linear-system}
    \begin{pmatrix} N^\top & 0 & 0 \\ H & -N & -\eta \\ \eta^\top & 0 & 0 \end{pmatrix}
    \begin{pmatrix} \beta \\ \lambda \\ \nu \end{pmatrix}
    = \begin{pmatrix} 0 \\ 0 \\ 1 \end{pmatrix},
    \qquad
    \beta \in \mathbb{R}^{|S|}, \quad \lambda \in \mathbb{R}^4, \quad \nu \in \mathbb{R}.
    % QC: KKT for max beta^T H beta subject to N^T beta = 0 (closure) and
    %   eta^T beta = 1 (normalization). Two constraints need two multipliers:
    %   lambda for closure, nu for normalization. Without nu, the system is
    %   overdetermined ((|S|+5) x (|S|+4)) and generically inconsistent.
  \end{equation}

\item \textbf{Filter.}
  Discard $(S, \sigma)$ if the system has no solution, or
  if $\beta_i \leq 0$ for any~$i$.
  If the system has multiple solutions, choose any solution
  with $\beta_i > 0$ for all~$i$, if one exists.
  % QC: the filter demands beta_i > 0 (strict), while M(K) in the theorem
  %   allows beta_i >= 0. This is intentional and correct:
  %   If beta_i = 0 for some i in S, that facet is unused and the same orbit
  %   appears as a solution for the strictly smaller subset S' = S \ {i}.
  %   So the maximum over M(K) is attained at an interior point (all beta_i > 0)
  %   of some sub-problem. The algorithm's enumeration over all subsets S
  %   catches this. No solutions are lost by the strict filter.

\item \textbf{Evaluate.}
  Record $\beta^\top H \beta$.
\end{enumerate}

\noindent
Return
\begin{equation}\label{eq:action-formula}
  c_{\mathrm{EHZ}}(K) = \bigl(\max_{S,\sigma} \beta^\top H \beta\bigr)^{-1}.
\end{equation}
\end{algorithm}

% Jörn: text approved (06a976c)
\begin{remark}\label{rem:critical-point}
The solution $\beta$ with $\beta_i > 0$
is the critical point of $\beta^\top H \beta$
subject to $N^\top \beta = 0$,
$\eta^\top \beta = 1$, and $\beta_i > 0$.
\end{remark}

% Jörn: text approved (06a976c) — from \begin{lemma} to \end{lemma}
\begin{lemma}[Well-definedness]\label{lem:well-defined}
% QC: proves that beta^T H beta = nu for any solution, so the value
%   recorded in Step~4 is independent of which solution the solver returns.
%   Key: use the system equations to replace H beta, then apply constraints.
If system~\eqref{eq:linear-system} has multiple solutions
$(\beta, \lambda, \nu)$, they all yield the same value
$\beta^\top H \beta$.
\end{lemma}

